% s05-sec3b.tex —— §3.3 Steinberg–Hitchin 基（p.25–30）、§3.4 椭圆正则元（p.30–33）
\subsection{Steinberg–Hitchin 纤维化的基}
\srcpage{25}

本文只处理迹公式的椭圆正则部分，它将用于形如 $f=\prod_vf_v$ 的函数，其中 $f_v$ 光滑紧支。为使用它，我们需要对稳定半单共轭类、也就是对我们所谓的 Steinberg–Hitchin 纤维化的基，有一个恰当的描述。我们先从半单、单连通且分裂的群开始，然后转向半单、单连通且拟分裂的群，接着是任意的半单单连通群，最后是我们所处理的最一般情形，即 z-扩张，我们提醒，此时导出群单连通。不过我们强调，稳定化隐含地意味着过渡到一个拟分裂群！

迹公式的主部是对椭圆正则共轭类求和；稳定迹公式的主部则是对稳定、椭圆、正则类求和。它们在环面情形与半单单连通群情形的描述不同，但在一般情形必须把两者混合起来，因为 $G$ 是由 $G_{\mathrm{der}}$ 与中心 $Z$ 构造的。除了一些尚待描述的上同调复杂性之外，$G(F)$ 中的半单共轭类（不妨说是稳定类）是 $Z(F)$ 中一个元素与 $G_{\mathrm{der}}(F)$ 中一个类的乘积。$G$ 上的局部或整体调和分析，归根结底归结为这两个因子上的调和分析。而 $G_{\mathrm{der}}(F)$ 中的稳定类与 $F$ 上有限维向量空间的点相联系，这就使我们可以用与以往完全不同的方式处理迹公式中出现的和。

对受过几何训练的读者，我们先把结论提前说出来，好让他们清楚迹公式带来的分析问题，以及我们所描述的加法结构或许提供的解决办法，也不至于被自己的几何知识束缚太多。我们尤其要提醒他们：域 $F$ 既可以是代数数域，也可以是函数域！

结论将是：稳定迹公式的椭圆部分，是对 $\eta\in h$（集合 $h$ 尚待描述）求和，再对集合 (3.15)（稍后引入）求和。

\srcpage{26}
若除以 $|A(F)|$，甚至可以就在 $B_\eta(F)\times Z_\eta(F)$ 上求和。尽管对 $\eta$ 的和不是有限的，但对给定的函数 $f=\prod_vf_v$，只有有限个 $\eta$ 的贡献不为 $0$。因此只需处理单个 $B_\eta(F)\times Z_\eta(F)$ 的贡献。群 $Z(F)$ 在 $Z_\eta(F)$ 上的作用是单传递的，所以对 $Z_\eta(F)$ 的求和就是对 $Z(F)$ 的求和，从而可以用关于偶 $Z(F)\subset Z(\mathbb A_F)$——更确切地说 $Z_+(F)\subset Z(\mathbb A_F)$——的 Poisson 公式处理。此处的新意，也是此前未出现过的，大致是：$B_\eta(F)$ 是向量空间，且对给定的 $z\in Z_\eta(F)$，对 $B_\eta(F)$ 的求和是某个 adèle 函数在 $B_\eta(F)$ 各点取值之和的 Poisson 和，而这个函数的行为足够好，可以使用 Poisson 公式。"大致"一词指的是引理 (4.1) 与 (4.2) 的推论。为使用 Poisson 公式需要截断，否则无法验证公式中出现的对偶和收敛。为此需要一些估计，而这些估计只能局部得到，而且即便如此，除少数将在 [L5] 中给出的特殊情形外也尚未得到。问题只在由引理 (4.1) 所允许的截断之后才出现。

这些问题只有在 Getz 引理所允许的截断之后才会出现；本文引入该引理，主要是为了能够陈述命题 5.6，但它的目标远不止于此。

在解释完所有定义之后，我们将给出 $\operatorname{GL}(2)$ 的简单例子，以清楚显示我们建议在 Poisson 公式中使用的迹公式的那一部分。

设 $T$ 是 $G$ 的一个环面，假定它分裂，且暂时假定 $G$ 半单单连通。设 $\alpha_1,\dots,\alpha_r$ 是 $T$ 关于任意选定序的单根。与这些单根相伴的是基本权 $\mu_1,\dots,\mu_r$，由 $\mu_i(\check\alpha_j)=\delta_{i,j}$ 定义。设 $\rho_i=\rho_{\mu_i}$ 是最高权为 $\mu_i$ 的表示，$b_i(t)=\operatorname{tr}\rho_i(t)$。诸 $b_i$ 在 $F$ 上代数独立，$T$ 对 Weyl 群的商是仿射空间 $\operatorname{Spec}F[b_1,\dots,b_r]$。这个商就是 Steinberg–Hitchin 基 $A$，其中 $b_1,\dots,b_r$ 将是我们偏爱的坐标。相反，$T$ 上的坐标取为 $\gamma_i=t^{\lambda_i}$，其中
\begin{equation}\tag{3.13}
\lambda_i=\sum_ja_{i,j}\mu_j,\qquad \det(a_{i,j})=\pm1.
\end{equation}
整数矩阵 $(a_{i,j})$ 的具体选取无关紧要。此时基 $A$ 也是它的线性部分，我们记作 $B$。当 $G\ne G_{\mathrm{der}}$ 时不再如此。

\srcpage{27}
Dynkin 图的自同构有限群提升为半单、单连通、分裂群 $G$ 的自同构群，其中每个自同构都定义在 $F$ 上。一个拟分裂但单连通半单的群由取值于这个自同构群的上闭链定义，也就是由同态 $\sigma\mapsto\varphi(\sigma):\operatorname{Gal}(\bar F^{\mathrm{sep}}/F)=\operatorname{Gal}_F\to\operatorname{Aut}(\text{Dynkin 图})$ 定义。后一个群也作用在 $T$（一个分裂 Cartan 子群）与 $G$ 上，这些作用都定义在 $F$ 上。

于是 $\operatorname{Gal}_F$ 在基本表示集合 $\rho_1,\dots,\rho_r$ 上有右作用
\[
\sigma:\rho_i\mapsto\rho_{i\sigma},\qquad
\rho_{i\sigma}(g)=\rho_i(\varphi(\sigma)(g)),
\]
从而在它们的特征 $b_1,\dots,b_r$ 上也有右作用。注意
\[
\rho_{(i\tau)\sigma}(g)=\rho_{i\tau}(\varphi(\sigma)(g))
=\rho_i(\varphi(\tau)\varphi(\sigma)(g))
=\rho_i(\varphi(\tau\sigma)(g))=\rho_{i\tau\sigma}(g).
\]
暂时设 $G_\varphi$ 是由上闭链 $\sigma\mapsto\varphi(\sigma)$ 定义的拟分裂群。Galois 群在 $G_\varphi(\bar F^{\mathrm{sep}})$ 上的作用（作为集合它只是 $G(\bar F^{\mathrm{sep}})$）为
\[
\sigma_\varphi:g\mapsto\varphi(\sigma)(\sigma(g))=\sigma(\varphi(\sigma)(g)),
\]
其中 $\sigma(g)$ 是相对于 $G$ 定义的。

若 $\{h_\sigma\}$ 是相对于这个作用、取值于 $G_\varphi(\bar F^{\mathrm{sep}})=G(\bar F^{\mathrm{sep}})$ 的中心的任一上闭链，则 $\rho_i(h_\sigma)\in\bar F^{\mathrm{sep}}$ 是一个纯量 $\epsilon_i(\sigma)$。此外 $\sigma(\epsilon_{i\sigma}(\tau))\epsilon_i(\sigma)=\epsilon_i(\sigma\tau)$，因为
\[
\sigma(\rho_{i\sigma}(h_\tau))\rho_i(h_\sigma)
=\rho_{i\sigma}(\sigma(h_\tau))\rho_i(h_\sigma)
=\rho_i(\varphi(\sigma)(\sigma(h_\tau))h_\sigma)=\rho_i(h_{\sigma\tau}).
\]
因此我们可以用上闭链
\begin{equation}\tag{3.14}
\sigma:(b_1,\dots,b_r)\mapsto(b'_1,\dots,b'_r),
\qquad b'_i=\epsilon_i(\sigma)b_{i\sigma}
\end{equation}
来扭转 $F$ 上以 $b_1,\dots,b_r$ 为坐标的线性空间。特别地，取 $\{h_\sigma\}$ 为平凡上闭链，我们就得到半单、单连通、分裂群的 Steinberg–Hitchin 基的一个扭形式，它同时是与同态 $\varphi$ 相联系的拟分裂群 $G_\varphi$ 的任何内形式的 Steinberg–Hitchin 基，特别是 $G_\varphi$ 本身的基。这个基始终是 $F$ 上的向量空间。

现在我们转向一般情形，即 $G$ 的导出群 $G_{\mathrm{der}}$ 单连通。为避免不可分问题，我们假定 $F$ 的特征与 (3.1) 中群 $A$ 的阶 $|A|$ 互素。于是我们又一次回避了一个小困难。

\srcpage{28}
按 [BT] 的命题 2.2，$G(\bar F^{\mathrm{sep}})$ 的每个元素 $g$——从而特别是 $G(F)$ 的每个元素——都是乘积 $g=g_1g_2$，其中 $g_1\in G_{\mathrm{der}}(\bar F^{\mathrm{sep}})$，$g_2\in Z(\bar F^{\mathrm{sep}})$。若 $g\in G(F)$，则对每个 $\sigma\in\operatorname{Gal}_F$ 有 $\sigma(g_1)^{-1}g_1\in A(\bar F^{\mathrm{sep}})$，且它等于 $\sigma(g_2)g_2^{-1}$。由于 $g_1$ 只在一个 $A(\bar F^{\mathrm{sep}})$ 的元素的意义下给定，只有上闭链 $\{\sigma(g_1)^{-1}g_1\}$ 在 $H^1(F,A)$ 中的像 $\eta(g)$ 是确定的。设 $h$ 是这样得到的 $H^1(F,A)$ 的子集。

我们想把 $G$ 的 Steinberg–Hitchin 基 $A=A_G$ 上的 $F$-有理点集，描述为对 $\eta\in h$——甚至对 $H^1(F,A)$——求和，求的是一个商的有理点集，即
\begin{equation}\tag{3.15}
\{B_\eta(F)\times Z_\eta(F)\}/A(F),
\end{equation}
其中 $B_\eta$ 是一个线性空间，是 $G_{\mathrm{der}}$ 的 Steinberg–Hitchin 基之下的线性空间的扭形式，$A(F)$ 作为 $G_{\mathrm{der}}(F)$ 的中心的子群作用在其上；$Z_\eta$ 是由 $\eta$ 定义的 $Z$ 上的 torsor，$A(F)$ 作为 $Z(F)$ 的子群作用在其上。为此，对每个 $\eta\in h$，选取并固定一个代表它的上闭链 $h=\{h_\sigma\}$。若 $\eta$ 或 $h$ 在 $H^1(F,Z)$ 中不变成平凡，则 $Z_\eta(F)$ 为空。反之，若它变成平凡，则 $Z_\eta(F)$ 可以与对这个给定的 $h$ 而言所有可能的 $g_2$ 的集合等同。使 $\{\sigma(z)z^{-1}\}=\{h_\sigma\}$ 的所有 $z\in Z(\bar F^{\mathrm{sep}})$ 构成 torsor $Z_\eta$ 的点集，也就是对这个 $h$ 而言所有可能的 $g_2$ 的集合。

假定
\begin{equation}\tag{3.16}
\sigma(g_1)^{-1}g_1=h_\sigma
\end{equation}
对一切 $\sigma$ 成立，并设 $b_1(g_1),\dots,b_r(g_1)$ 是 $g_1$ 在 $G_{\mathrm{der}}$ 的 Steinberg–Hitchin 基 $A_{\mathrm{der}}$ 中的像 $\beta(g_1)$ 的线性坐标。我们已指出，对拟分裂半单单连通群的任何内扭形式，这个基作为 $F$ 上的簇，等于我们所描述其结构的拟分裂群的 Steinberg–Hitchin 基。特别地，它是线性空间，是分裂群之基的扭形式。我们按等式 (3.14) 用与 $\{h_\sigma\}$ 相联系的上闭链扭转这个 $F$ 上维数为 $r$ 的线性空间。我们验证条件 (3.16) 等价于 $\beta(g_1)$ 在扭空间中是有理的。我们强调，此时基本的线性空间是拟分裂群的 Steinberg–Hitchin 基！

按 (3.14)，对一切 $i$ 的有理性条件 $b'_i=b_i$ 变成
\[
b'_i(\beta(g_1))=b'_i(g_1)=\epsilon_i(\sigma)\sigma(b_{i\sigma}(g_1)).
\]

\srcpage{29}
有 $\beta(\sigma(g_1))=\sigma(\beta(g_1))$。更确切地说，由于 $G_{\mathrm{der}}$ 的 Steinberg–Hitchin 基是已描述的半单、单连通、分裂群（$G_{\mathrm{der}}$ 本身是其扭形式）的 Steinberg–Hitchin 基的扭形式，有
\[
b_i(\sigma(g_1))=\sigma(b_{i\sigma}(g_1)).
\]
也就是说，$\operatorname{Gal}_F$ 在基上的作用由 (3.14) 给出，上闭链 $\{\epsilon_i(\sigma)\}$ 平凡，故
\[
\sigma:(b_i)\mapsto(\sigma(b_{i\sigma})).
\]
尽管 Galois 群在 Steinberg–Hitchin 基上的作用——更确切地说，在借助坐标 $b_i$ 与之相联系的线性空间上的作用——是它在分裂群之基上的作用，但它在 $G_{\mathrm{der}}(\bar F^{\mathrm{sep}})$ 上的作用却是由它在 $G(\bar F^{\mathrm{sep}})$ 上的作用导出的，而不是由它在分裂群或拟分裂群上的作用导出的。

我们承认，把这些定义都塞进文章的行李中是乏味的；但半单单连通群的 Steinberg–Hitchin 基的线性结构，对我们的目的而言是迹公式的基本要素之一。

在处理 3.1 节所描述的那类约化群时，必须引入 (3.16) 的上闭链 $\{\epsilon_i(\sigma)\}$，以及由 (3.14) 与之相联系的分裂半单单连通群之基的扭形式。我们断言，相对于这个形式，像 $\beta(g_1)$ 定义在 $F$ 上，因此
\[
b_i(\beta(g_1))=b_i(g_1)=\epsilon_i(\sigma)\sigma(b_{i\sigma}(g_1)).
\]
事实上，
\[
\sigma(b_{i\sigma}(g_1))=b_i(\sigma(g_1))=b_i(h_\sigma^{-1}g_1)=\epsilon_i^{-1}(\sigma)b_i(g_1).
\]

我们无法回避由 $G$ 的中心中的分裂环面带来的麻烦。此时要处理的共轭类不是 $G(F)$ 的类，而是 $G'(F)$ 的类。由于每个这样的类都含有 $\tilde Z$ 中一个元素与 $G(F)$ 中一个元素的乘积 $zt$，描述 (3.15) 大致仍然正确。只需把 $Z_\eta$ 换成离散群 $Z^+_\eta=Z_{s_G}Z_\eta$，因为若 $z't'=g^{-1}ztg$，其中 $t',t,g\in G(F)$，$z,z'\in Z_{s_G}$，则
\[
z^{-1}z'\in Z_{s_G}\cap G(F)=\{1\}.
\]
我们记 $A(F)$ 为 $G$ 的 Steinberg–Hitchin 基上的 $F$-点集，对 $F_v$ 或 $\mathbb A_F$ 用类似记号。设 $c(g)$ 是 $g\in G$ 在 $A$ 中的像。按刚才的验证，它是 $B_\eta(F)\times Z^+_\eta(F)$ 模 $A(F)$ 之商对 $\eta$ 的并。集合 $Z_\eta$ 是 $F$ 上的 torsor，从而是 $Z^+(F)$ 的主齐性空间，嵌入 $Z(\mathbb A_F)$。另一方面，$B_\eta(F)$ 是

\srcpage{30}
有限维向量空间，嵌入 $B_\eta(\mathbb A_F)$，后者是 $\mathbb A_F$ 上有限秩自由模。于是为利用迹公式而想到的技术，就是对 $Z^+(F)\subset Z(\mathbb A_F)$ 与 $B_\eta(F)\subset B_\eta(\mathbb A_F)$ 使用 Poisson 公式。$Z$ 与 $B_\eta$ 的整体与局部测度按 3.2 节选取。

所有这些解释做完之后，显然基本情形是半单单连通群的情形，本文最后几节将只限于这一情形。不过，对 $\operatorname{GL}(2)$ 补充几点精确说明仍是有用的。Steinberg–Hitchin 基是二维的，映射 $c$ 由 $t\mapsto(b,a)$ 给出，其中 $X^2-bX+a$ 是 $t$ 的特征多项式。参数 $\eta$ 是 $a$ 模 $F^\times$ 的平方元集合 $c^2$ 的一个选取 $a_0$，也就是上闭链 $\sigma\mapsto\sigma(\sqrt{a_0})/\sqrt{a_0}$ 的一个选取。于是写作 $a=c^2a_0$，$b=d\sqrt{a_0}$，使得 $d=b/\sqrt{a_0}$ 是加法群上的一个主齐性空间中坐标为 $F$ 中元素的点。对 $b$——或对 $d$——的求和将使用加法 Poisson 对偶。这一描述即使在特征 $2$ 的域上也成立，但此时 $F$ 上合适的拓扑不是平展拓扑。这些是留待另一次讨论的问题。

\subsection{稳定迹公式中的椭圆正则元}
\srcpage{30}

本文这一部分，参考文献中 Kottwitz 与 Ono 的工作是根本性的。\footnote{由于我们原则上不仅想处理 $F$ 为数域的情形，也想处理 $F$ 为函数域的情形，这些文献并不完全够用。不过我们决定暂时不处理这一点。} 笔记 [L2] 也可能有用。我们始终假定 $G_{\mathrm{der}}$ 拟分裂且单连通。

我们提醒，为得到稳定公式的椭圆正则部分，我们从和
\begin{equation}\tag{3.17}
\sum_\gamma\meas^{\mathrm{prod}}\bigl(T^+_\gamma(F)\backslash T_\gamma(\mathbb A_F)\bigr)
\int_{T_\gamma(\mathbb A_F)\backslash G(\mathbb A_F)}f(g^{-1}\gamma g)\,d\bar g
\end{equation}
开始。第一项的测度由 $d^{\mathrm{prod}}t$ 给出。第二项的测度 $d\bar g$ 是商测度
\[
\frac{d^{\mathrm{prod}}g}{d^{\mathrm{prod}}t}=\prod_vd\bar g_v,
\qquad
\bar g_v=\frac{L(1,\sigma_G)}{L(1,\sigma_T)}\,\frac{d^{\mathrm{geom}}g_v}{d^{\mathrm{geom}}t_v}.
\]

\srcpage{31}
公式 (3.17) 的形式验证很容易。我们算子的核是
\[
\sum_{\gamma\in T^+(F)}f(g^{-1}\gamma h).
\]
迹形式上由积分
\[
\int_{G_+(F)\backslash G(\mathbb A_F)}f(g^{-1}\gamma g)\,dg
\]
得到，它形式上等于对 $G_+(F)$ 中共轭类 $\{\gamma\}$ 的和
\begin{equation}\tag{3.18}
\int_{G^+_\gamma(F)\backslash G(\mathbb A_F)}f(g^{-1}\gamma g)\,dg.
\end{equation}
事实上这个和一般没有意义。不过，在真正的迹公式中（例如见 [A2]），出现的是对 $G_+(F)$ 中椭圆正则共轭类求和所得的各表达式之和。本文感兴趣的是这个和，因为一切迹公式的研究都从椭圆正则类的贡献开始。若 $\gamma$ 是椭圆正则的，则其中心化子是环面，从而连通。\footnote{我们暂时不给出参考文献。还需要在正特征下为 Kottwitz 的定理 4.4（[K1]）找一个证明。}

既然是稳定迹公式，我们处理的就是稳定共轭类而非共轭类。按 (3.15)，这些类的参数是 Steinberg–Hitchin 基 $B_\eta(F)$ 中的点。若 $G$ 不等于其导出群，则 $Z_{s_G}$ 与 $Z_\eta(F)$ 的存在使公式更复杂。正因如此我们假定了 $G=G_{\mathrm{der}}$；但为了不时提醒自己这不是最一般的情形，我们有时使用兼顾一般情形的记号。重要的是 $B_\eta$ 是 $F$ 上的向量空间。

我们考察 (3.18) 中的项，并加上对 $Z_{s_G}$ 的求和，以免忘记起作用的是 $G_+(F)$ 中的共轭类而非 $G(F)$ 中的，于是考虑
\begin{equation}\tag{3.19}
\sum_{z\in Z_{s_G}}\int_{G^+_\gamma(F)\backslash G(\mathbb A_F)}f(g^{-1}z\gamma g)\,dg.
\end{equation}
群 $G_\gamma$（有时记作 $T_\gamma$）是正则半单元 $\gamma$ 的中心化子。在 [L2] 中，迹公式的正则椭圆部分被转化为对稳定迹公式的内窥群求和。我们在此只关心

\srcpage{32}
$G$ 自身的贡献。我们提醒，$G$ 始终满足 $G_{\mathrm{der}}$ 单连通，且此刻为简化起见等于 $G_{\mathrm{der}}$。我们尽量引用 [K2] 简要说明它的形式，因为那里的记号与解释比 [L2] 更易理解。

首先，(3.19) 中的积分
\[
\int_{G^+_\gamma(F)\backslash G(\mathbb A_F)}f(g^{-1}z\gamma g)\,d^{\mathrm{prod}}g
\]
写成如下形式：
\begin{equation}\tag{3.20}
\meas^{\mathrm{prod}}\bigl(T^+_\gamma(F)\backslash T_\gamma(\mathbb A_F)\bigr)
\int_{T(\mathbb A_F)\backslash G(\mathbb A_F)}f(g^{-1}\gamma g)\,d^{\mathrm{prod}}\bar g.
\end{equation}
作了这个几乎是形式的修改之后，先不对 $G(F)$ 的所有正则类 $\{\gamma\}$ 求和（如在 (3.3) 中那样），而对在 $G(\mathbb A_F)$ 中共轭于一个给定 $\gamma$ 的所有 $\gamma$ 求和。\yizhu{原文此处引 (3.3) 与 (3.1)，对照上下文，所指当为 (3.19) 与 (3.17)；这两处交叉引用编号有误，译文照录原文编号。}{依据：原书 p.32；对照 (3.17)–(3.19)} 这给出一个因子，在 [K2] 中记作
\[
\bigl|\ker[E(T/F)\to E(T/\mathbb A_F)]\bigr|,
\]
其中 $T=T_\gamma$。但随后为了得到对内窥群的和，要引入关于群 $K(T/F)$ 的特征的展开。至少当只考虑内窥群 $G$ 本身、也就是稳定部分的贡献时，这导致 [K2] 中的系数
\begin{equation}\tag{3.21}
\iota(F,T,G)=\frac{\bigl|\ker[E(T/F)\to E(T/\mathbb A_F)]\bigr|}{K(T/F)},
\end{equation}
记号与 [L2] 完全相同。在 [K2] 的引理 8.3.2 中——该文写于关于 Tamagawa 数的 Weil 猜想（[K3]）与关于 Hasse 原理的 Kneser 猜想的完整证明之前——验证了
\begin{equation}\tag{3.22}
\iota(F,T,G)=\frac{\tau(G)}{\tau(T)}.
\end{equation}
引入因子 (3.21) 并过渡到稳定迹公式，其后果是 (3.1) 中对 $\gamma$ 的和被对稳定椭圆正则类的和代替，从而实际上是对 Steinberg–Hitchin 基中椭圆正则元的和。对椭圆 Cartan 子群，$m_T=m_G$，且由 (3.10)，(3.20) 中的因子
\[
\meas^{\mathrm{prod}}\bigl(T^+_\gamma(F)\backslash T_\gamma(\mathbb A_F)\bigr)
=m_T\rho_T\tau(T)=m_G\rho_T\tau(T)
\]

\srcpage{33}
乘以 (3.22) 等于
\begin{equation}\tag{3.23}
m_G\tau(G)\rho_T=m_G\rho_G\tau(G)\frac{\rho_T}{\rho_G}.
\end{equation}
剩下的是
\[
\int_{T(\mathbb A_F)\backslash G(\mathbb A_F)}f(g^{-1}\gamma g)\,d^{\mathrm{prod}}\bar g
=\prod_v\int_{T(F_v)\backslash G(F_v)}f_v(g^{-1}\gamma g)\,d\bar g_v.
\]
我们引入记号
\[
\int_{T(F_v)\backslash G(F_v)}f_v(g^{-1}\gamma g)\,d\bar g_v=\Orb(\gamma,f_v).
\]
把这些最后的积分对 $\gamma$ 的稳定共轭类中的共轭类代表求和，就得到稳定轨道积分 $\Orb(\gamma_{\mathrm{st}},f_v)$，其中 $\gamma_{\mathrm{st}}$ 记 $\gamma$ 的稳定类。在迹公式中出现的是乘积
\begin{equation}\tag{3.24}
\prod_v\Orb(\gamma_{\mathrm{st}},f_v).
\end{equation}
若用精确但笨拙的记号，就得把下标 "st" 写两次，再加到 "Orb" 上。

因子 $m_G\rho_G\tau(G)$ 与 $T$ 无关。因此我们暂时把它搁置一旁。其中隐藏着一些上同调与测度的结果，我们以后不再显式需要它们。表示 $\sigma_G$ 是 $\sigma_T$ 的子表示，设 $\sigma_{T/G}$ 为商。对椭圆 $T$，这个商不含平凡表示，且
\[
\frac{\rho_T}{\rho_G}=L(1,\sigma_{T/G})=\lim_{s\downarrow1}L(s,\sigma_{T/G}).
\]
下一部分我们感兴趣的是乘积
\begin{equation}\tag{3.25}
L(s,\sigma_{T/G})\prod_v\Orb(\gamma_{\mathrm{st}},f_v),
\qquad s>1.
\end{equation}
这个乘积对 $T(F)$ 中任何稳定共轭类都有定义，即使它不是椭圆的。然而当 $s\to1$ 时，若该类不是椭圆的，极限不存在。由于我们从不只使用一个迹公式，而总是使用若干个带不同符号的迹公式之和，所以完全可以把在某一个公式中取极限时发散的部分加进来，只要这部分在另一个公式中以相反符号出现。本文中不会有例子，但后续文章会有。
